% !TEX root = ../main.tex
%% Corrigé du sujet type bac n° 1 — valeurs vérifiées : verifications/sujets_bac.py

\corrigeSujet{1}
{\otLight\fontsize{8.5}{11.5}\selectfont\color{otMuted}\noindent Les barèmes répartissent les 20 points de chaque sujet entre les
questions. Dans une correction officielle, deux de ces points évaluent la maîtrise de la langue sur l'ensemble de la copie :
orthographe, syntaxe, rigueur du raisonnement et précision du vocabulaire (grille du guide de l'épreuve).\par}

\begin{corrigeBac}{9}{Un détartrant à l'acide lactique}
\partieExo{1}{La molécule d'acide lactique}
\begin{questions}
\item Le groupe \res{\text{hydroxyle}} $-\mathrm{OH}$, porté par un carbone tétragonal, caractérise la famille des
\res{\text{alcools}} ; le groupe \res{\text{carboxyle}} $-\mathrm{COOH}$ celle des \res{\text{acides carboxyliques}}. \pts{0,75 pt}
\item Une bande forte et large vers 3200 -- \qty{3400}{cm^{-1}} (O--H de l'alcool), une bande très large de 2500 à
\qty{3200}{cm^{-1}} (O--H de l'acide) et une bande forte et fine vers 1700 -- \qty{1725}{cm^{-1}} (C$=$O). \pts{0,5 pt}
\item $\ce{AH + H2O <=> A- + H3O+}$ : l'acide lactique cède un proton à l'eau. \pts{0,5 pt}
\end{questions}
\partieExo{2}{Contrôler l'étiquette d'un détartrant}
\begin{questions}[resume]
\item Un litre de détartrant a pour masse \qty{1100}{g} et contient $0{,}45\times1100=\qty{495}{g}$ d'acide lactique, soit
$c=\dfrac{495}{90{,}1}=\res{\qty{5,5}{mol.L^{-1}}}$. La solution S est 100 fois moins concentrée, et à l'équivalence
$c_S\,V_A=c_B\,V_\mathrm{E}$ : $V_\mathrm{E}=\dfrac{0{,}0549\times20{,}0}{0{,}100}=\res{\qty{11,0}{mL}}$. \pts{1 pt}
\item On prélève \qty{1,0}{mL} de détartrant avec la \res{\text{pipette jaugée de 1,0 mL}}, que l'on verse dans la
\res{\text{fiole jaugée de 100,0 mL}} ; on complète avec de l'eau distillée jusqu'au trait de jauge, puis on homogénéise.
La verrerie jaugée garantit un facteur de dilution précis : $100{,}0/1{,}0=100$. \pts{0,5 pt}
\item $\ce{AH + HO- -> A- + H2O}$. \pts{0,5 pt}
\item Méthode des tangentes : on trace deux tangentes parallèles de part et d'autre du saut de pH, puis la parallèle
équidistante, qui coupe la courbe au point d'équivalence : $\res{V_\mathrm{E}=\qty{10,9}{mL}}$ (pH voisin de 8,2). \pts{1 pt}
\item D'après le titrage, S a une concentration voisine de \qty{0,0545}{mol.L^{-1}} ; un acide fort donnerait
$\pH=-\log0{,}0545=\res{1{,}3}$. La courbe indique un pH initial de 2,6 : la réaction de l'acide avec l'eau n'est pas totale,
l'acide lactique est un \res{\text{acide faible}}. \pts{0,75 pt}
\item À la demi-équivalence, $V_B=V_\mathrm{E}/2=\qty{5,45}{mL}$, la moitié de l'acide a été convertie : $[\ce{AH}]=[\ce{A-}]$, donc
$\pH=\pKa$. La courbe donne $\res{\pKa\approx3{,}9}$, en accord avec la donnée. \pts{0,5 pt}
\item Le pH à l'équivalence, voisin de 8,2, est dans le saut de pH (de 6 à 11 environ) : la \res{\text{phénolphtaléine}}, dont la
zone de virage est contenue dans le saut, convient. \pts{0,5 pt}
\item $c_S=\dfrac{c_B\,V_\mathrm{E}}{V_A}=\dfrac{0{,}100\times10{,}9}{20{,}0}=\qty{0,0545}{mol.L^{-1}}$, donc $c=100\,c_S=\qty{5,45}{mol.L^{-1}}$. Un litre de détartrant
contient $5{,}45\times90{,}1=\qty{491}{g}$ d'acide pour une masse de \qty{1100}{g} : titre massique $t=\dfrac{491}{1100}=\res{44{,}6~\%}$.
Son incertitude-type vaut $u(t)=0{,}011\times44{,}6=\qty{0,49}{\percent}$, et l'écart normalisé
$z=\dfrac{|44{,}6-45|}{0{,}49}=0{,}8<2$ : le titre mesuré est \res{\text{compatible}} avec l'étiquette. \pts{2 pts}
\end{questions}
\partieExo{3}{L'acide lactique dans le muscle}
\begin{questions}[resume]
\item \ce{AH} prédomine pour $\pH<3{,}86$, \ce{A-} pour $\pH>3{,}86$. À $\pH=7{,}4$, c'est l'\res{\text{ion lactate}} qui prédomine. \pts{0,5 pt}
\item $\dfrac{[\ce{A-}]}{[\ce{AH}]}=10^{\pH-\pKa}=10^{7{,}4-3{,}86}=\res{\num{3.5e3}}$ : l'acide lactique est presque entièrement sous forme de lactate
dans la cellule. \pts{0,5 pt}
\end{questions}
\schema[corrige]{bac-s1-ph}
\end{corrigeBac}

\begin{corrigeBac}{6}{Le lancer du poids}
\begin{questions}
\item $P=m\,g=7{,}26\times9{,}81=\res{\qty{71,2}{N}}$. Avec $S=\pi\times0{,}060^2=\qty{1,13e-2}{m^2}$,
$F=\frac12\times1{,}2\times0{,}47\times\num{1.13e-2}\times13{,}5^2=\res{\qty{0,58}{N}}$, moins de 1~\% du poids : les frottements sont
négligeables. \pts{0,5 pt}
\item La boule n'est soumise qu'à son poids : $m\,\vec a=m\,\vec g$, soit $\vec a=\vec g$, de coordonnées $\res{a_x=0}$ et
$\res{a_y=-g}$. \pts{0,75 pt}
\item En intégrant avec $v_x(0)=v_0\cos\alpha=\qty{10,3}{m.s^{-1}}$ et $v_y(0)=v_0\sin\alpha=\qty{8,68}{m.s^{-1}}$ : $v_x=v_0\cos\alpha$ et
$v_y=-g\,t+v_0\sin\alpha$. Puis, avec $x(0)=0$ et $y(0)=h$ : $\res{x=(v_0\cos\alpha)\,t}$ et
$\res{y=-\frac12\,g\,t^2+(v_0\sin\alpha)\,t+h}$. \pts{0,75 pt}
\item $t=\dfrac{x}{v_0\cos\alpha}$, d'où $y=-\dfrac{g}{2\,v_0^2\cos^2\alpha}\,x^2+x\tan\alpha+h$, avec $\dfrac{9{,}81}{2\times10{,}34^2}=0{,}0459$ et
$\tan\ang{40}=0{,}839$ : $\res{y=-0{,}0459\,x^2+0{,}839\,x+2{,}10}$. \pts{0,75 pt}
\item Au sommet, $v_y=0$ : $t=\dfrac{v_0\sin\alpha}{g}=\qty{0,885}{s}$ et
$y_\mathrm{max}=h+\dfrac{(v_0\sin\alpha)^2}{2g}=2{,}10+\dfrac{8{,}68^2}{2\times9{,}81}=\res{\qty{5,94}{m}}$. \pts{0,75 pt}
\item L'impact correspond à $y=0$ : $-0{,}0459\,x^2+0{,}839\,x+2{,}10=0$. Le discriminant vaut
$0{,}839^2+4\times0{,}0459\times2{,}10=1{,}089$ et la racine positive
$x=\dfrac{0{,}839+\sqrt{1{,}089}}{2\times0{,}0459}=\res{\qty{20,5}{m}}$. \pts{1 pt}
\item Sans frottement, l'énergie mécanique se conserve : $\frac12\,m\,v^2=\frac12\,m\,v_0^2+m\,g\,h$, d'où
$v=\sqrt{13{,}5^2+2\times9{,}81\times2{,}10}=\res{\qty{14,9}{m.s^{-1}}}$. \pts{0,75 pt}
\item \code{t} est la durée du vol, racine positive de $y(t)=0$. Le programme calcule la longueur du jet pour chaque angle
entier de \ang{20} à \ang{60} et garde la plus grande : l'angle optimal vaut \res{\ang{42}}, pour \qty{20,57}{m}. La boule part
au-dessus du sol : la durée du vol est allongée, et il vaut mieux privilégier la vitesse horizontale, donc un angle
inférieur à \ang{45}. \pts{0,75 pt}
\end{questions}
\end{corrigeBac}

\begin{corrigeBac}{5}{Un capteur d'humidité}
\begin{questions}
\item Loi des mailles : $E=R\,i+u_C$, avec $i=C\,\dfrac{\mathrm{d}u_C}{\mathrm{d}t}$ : $\res{RC\,\dfrac{\mathrm{d}u_C}{\mathrm{d}t}+u_C=E}$. \pts{1 pt}
\item $\dfrac{\mathrm{d}u_C}{\mathrm{d}t}=\dfrac{E}{\tau}\,\mathrm{e}^{-t/\tau}$, d'où $RC\,\dfrac{E}{\tau}\,\mathrm{e}^{-t/\tau}+E-E\,\mathrm{e}^{-t/\tau}=E$ puisque $\tau=RC$ :
l'équation est vérifiée ; de plus $u_C(0)=E(1-1)=0$. \pts{0,75 pt}
\item À $t=\tau$, $u_C=0{,}63\,E=\qty{2,1}{V}$ : la courbe donne $\res{\tau=\qty{186}{\micro s}}$ (la tangente à l'origine coupe
l'asymptote $u_C=E$ à la même date). $C=\dfrac{\tau}{R}=\dfrac{\num{186e-6}}{\num{1.0e6}}=\res{\qty{186}{pF}}$. \pts{1 pt}
\item $\mathrm{HR}=\dfrac{186-161}{0{,}34}=\res{74~\%}$. \pts{0,75 pt}
\item À 70~\%, $C=161+0{,}34\times70=\qty{184,8}{pF}$, d'où $\tau=RC=\res{\qty{185}{\micro s}}$ : la ventilation se déclenche dès que
$\tau$ dépasse cette valeur. Un écart de \qty{1}{\micro s} correspond à 2,9~\% d'humidité : la durée doit être mesurée
à la microseconde près. \pts{1 pt}
\item $10\,\tau=\qty{1,86}{ms}$, d'où $\dfrac{1}{\num{1.86e-3}}=538$ : environ \res{\text{500 mesures par seconde}}. \pts{0,5 pt}
\end{questions}
\schema[corrige]{bac-s1-rc}
\end{corrigeBac}
